\textbf{Phần 1. Trắc nghiệm lựa chọn}
\setcounterpageref{ex}{0}

\begin{ex} %[2D4V22]
Kết quả của phép trừ $52 - 7$ là
\choice
{$43$}
{$44$}
{$46$}
{\True $45$}
\loigiai{
$52 - 7 = 45$.
\textbf{Đáp án của câu hỏi là D}
}
\end{ex}

\begin{ex} %[2D4V22]
Một đàn gà có $43$ con gà. Người ta bán đi $8$ con gà. Hỏi đàn gà còn lại bao nhiêu con gà?
\choice
{$34$ con gà}
{\True $35$ con gà}
{$36$ con gà}
{$33$ con gà}
\loigiai{
$43 - 8 = 35$ (con gà).
\textbf{Đáp án của câu hỏi là B}
}
\end{ex}

\begin{ex} %[2D4V22]
Tìm $x$ biết: $61 - x = 9$.
\choice
{$51$}
{$53$}
{\True $52$}
{$50$}
\loigiai{
$x = 61 - 9 = 52$.
\textbf{Đáp án của câu hỏi là C}
}
\end{ex}

\begin{ex} %[2D4V22]
Tính giá trị của biểu thức: $73 - 6 - 8$.
\choice
{\True $59$}
{$58$}
{$60$}
{$57$}
\loigiai{
$73 - 6 = 67$; $67 - 8 = 59$.
\textbf{Đáp án của câu hỏi là A}
}
\end{ex}

\begin{ex} %[2D4V22]
Một sợi dây dài $84\text{ cm}$, người ta cắt bớt đi $9\text{ cm}$. Chiều dài đoạn dây còn lại là
\choice
{$73\text{ cm}$}
{$74\text{ cm}$}
{$76\text{ cm}$}
{\True $75\text{ cm}$}
\loigiai{
$84 - 9 = 75\text{ (cm)}$.
\textbf{Đáp án của câu hỏi là D}
}
\end{ex}

\begin{ex} %[2D4V22]
Điền số thích hợp: $35 - 8 = 19 + \dots$
\choice
{$7$}
{\True $8$}
{$9$}
{$6$}
\loigiai{
Vế trái: $35 - 8 = 27$. Ta có: $27 - 19 = 8$. Vậy điền 8.
\textbf{Đáp án của câu hỏi là B}
}
\end{ex}

\begin{ex} %[2D4V22]
Năm nay mẹ $41$ tuổi, mẹ hơn con $8$ lần tuổi con là không phải, mẹ hơn con $32$ tuổi. Tuổi của con là
\choice
{$8$ tuổi}
{$10$ tuổi}
{\True $9$ tuổi}
{$7$ tuổi}
\loigiai{
$41 - 32 = 9$ (tuổi).
\textbf{Đáp án của câu hỏi là C}
}
\end{ex}

\begin{ex} %[2D4V22]
Cho phép tính: $\overline{a2} - 7 = 35$. Giá trị của chữ số $a$ là
\choice
{\True $4$}
{$3$}
{$5$}
{$2$}
\loigiai{
Ta có: $\overline{a2} = 35 + 7 = 42$. Do đó $a = 4$.
\textbf{Đáp án của câu hỏi là A}
}
\end{ex}

\begin{ex} %[2D4V22]
Một xe chở $91\text{ kg}$ gạo, đã dỡ xuống $8\text{ kg}$. Khối lượng gạo còn trên xe là
\choice
{$81\text{ kg}$}
{$82\text{ kg}$}
{$84\text{ kg}$}
{\True $83\text{ kg}$}
\loigiai{
$91 - 8 = 83\text{ (kg)}$.
\textbf{Đáp án của câu hỏi là D}
}
\end{ex}

\begin{ex} %[2D4V22]
Hiệu giữa số tròn chục liền sau số $60$ và số lớn nhất có một chữ số là
\choice
{$59$}
{\True $61$}
{$60$}
{$62$}
\loigiai{
Số tròn chục liền sau 60 là 70. Số lớn nhất có 1 chữ số là 9.\nHiệu: $70 - 9 = 61$.
\textbf{Đáp án của câu hỏi là B}
}
\end{ex}

\begin{ex} %[2D4V22]
Tìm một số biết rằng lấy số đó cộng với $6$ thì được $82$.
\choice
{$74$}
{$75$}
{\True $76$}
{$78$}
\loigiai{
Số đó là: $82 - 6 = 76$.
\textbf{Đáp án của câu hỏi là C}
}
\end{ex}

\begin{ex} %[2D4V22]
Điền dấu thích hợp: $63 - 5 \dots 49 + 8$.
\choice
{\True $>$}
{$<$}
{$=$}
{Không so sánh được}
\loigiai{
$63 - 5 = 58$; $49 + 8 = 57$. Vì $58 > 57$ nên điền dấu $>$. 
\textbf{Đáp án của câu hỏi là A}
}
\end{ex}

\textbf{Phần 2. Trắc nghiệm đúng sai}
\setcounterpageref{ex}{0}

\begin{ex} %[2D4V22]
Xét tính đúng sai của phép trừ có nhớ số có hai chữ số cho số có một chữ số:
\choiceTF
{\True Phép trừ $71 - 8$ có kết quả bằng $63$.}
{Trong phép trừ $54 - 6$, vì 4 không trừ được 6 nên ta lấy 6 trừ 4 bằng 2.}
{\True Biểu thức $82 - 9$ có giá trị bằng $73$.}
{\True Nếu số bị trừ giảm đi 4 đơn vị và số trừ giữ nguyên thì hiệu giảm đi 4 đơn vị.}
\loigiai{
\begin{itemize}
    \item Ý a) Đúng: $71 - 8 = 63$.
    \item Ý b) Sai: Sai lầm phổ biến của học sinh: phải mượn 1 chục thành 14 trừ 6 bằng 8.
    \item Ý c) Đúng: $82 - 9 = 73$. Đúng.
    \item Ý d) Đúng: Đúng theo tính chất phép trừ.
\end{itemize}
\begin{itemize}
    \item a) \textbf{Đúng}.
    \item b) \textbf{Sai}.
    \item c) \textbf{Đúng}.
    \item d) \textbf{Đúng}.
\end{itemize}
}
\end{ex}

\begin{ex} %[2D4V22]
Bác Năm thu hoạch được 62 quả dừa, bác đã bán đi 9 quả dừa:
\choiceTF
{\True Bác Năm còn lại 53 quả dừa.}
{\True Số dừa đã bán ít hơn số dừa còn lại là 44 quả.}
{Nếu bác bán thêm 3 quả nữa thì bác còn lại 51 quả.}
{Số 53 là số chẵn.}
\loigiai{
\begin{itemize}
    \item Ý a) Đúng: $62 - 9 = 53$ quả.
    \item Ý b) Đúng: $53 - 9 = 44$ quả.
    \item Ý c) Sai: Số dừa lúc này còn: $53 - 3 = 50$ quả (không phải 51 quả).
    \item Ý d) Sai: 53 có chữ số tận cùng là 3 nên là số lẻ.
\end{itemize}
\begin{itemize}
    \item a) \textbf{Đúng}.
    \item b) \textbf{Đúng}.
    \item c) \textbf{Sai}.
    \item d) \textbf{Sai}.
\end{itemize}
}
\end{ex}

\textbf{Phần 3. Trả lời ngắn}
\setcounterpageref{ex}{0}

\begin{ex} %[2D4V22]
Tính kết quả: $75 - 8$.
\shortans{67}
\loigiai{
\textbf{Đáp số:} $67$\\
$75 - 8 = 67$.
}
\end{ex}

\begin{ex} %[2D4V22]
Tìm số tự nhiên x biết: $x + 9 = 62$.
\shortans{53}
\loigiai{
\textbf{Đáp số:} $53$\\
$x = 62 - 9 = 53$.
}
\end{ex}

\begin{ex} %[2D4V22]
Một giá sách có 51 quyển sách. Người ta mượn đi 6 quyển sách. Hỏi trên giá còn lại bao nhiêu quyển sách?
\shortans{45}
\loigiai{
\textbf{Đáp số:} $45$\\
$51 - 6 = 45$ quyển.
}
\end{ex}

\begin{ex} %[2D4V22]
Một sợi dây dài 83 cm, người ta cắt đi 7 cm. Hỏi sợi dây còn lại dài bao nhiêu xăng-ti-mét?
\shortans{76}
\loigiai{
\textbf{Đáp số:} $76$\\
$83 - 7 = 76\text{ cm}$.
}
\end{ex}

\begin{ex} %[2D4V22]
Tìm hiệu của số bé nhất có hai chữ số mà chữ số hàng chục là 9 và số 8.
\shortans{82}
\loigiai{
\textbf{Đáp số:} $82$\\
Số bé nhất có 2 chữ số hàng chục là 9 là 90. Hiệu: $90 - 8 = 82$.
}
\end{ex}

\begin{ex} %[2D4V22]
Lớp 2A có 34 học sinh, trong đó có 8 bạn học sinh giỏi. Hỏi lớp 2A có bao nhiêu bạn không đạt học sinh giỏi?
\shortans{26}
\loigiai{
\textbf{Đáp số:} $26$\\
$34 - 8 = 26$ bạn.
}
\end{ex}
